% Generated draft: requires mathematical and source review.
\documentclass[11pt]{article}
\usepackage[margin=25mm]{geometry}
\usepackage{fontspec}
\setmainfont{Latin Modern Roman}
\usepackage{amsmath,amssymb,mathtools,unicode-math}
\setmathfont{Latin Modern Math}
\usepackage{enumitem,needspace,xurl,hyperref}
\setlength{\parindent}{0pt}
\setlength{\parskip}{5pt}
\setlength{\emergencystretch}{3em}
\newcommand{\N}{\mathbb{N}}
\newcommand{\Z}{\mathbb{Z}}
\newcommand{\Q}{\mathbb{Q}}
\newcommand{\R}{\mathbb{R}}
\newcommand{\C}{\mathbb{C}}
\newcommand{\F}{\mathbb{F}}
\newcommand{\A}{\mathbb{A}}
\newcommand{\PP}{\mathbb{P}}
\newcommand{\E}{\mathbb{E}}
\newcommand{\eps}{\varepsilon}
\newcommand{\dd}{\,\mathrm{d}}
\begin{document}
\section*{UnsolvedMath ID: 30000326}
\textbf{Source problem code:} catalogue entry
\subsection{Definitions and mathematical statement}
\textbf{Proposed self-intersection formula, as supplied.} For a regular arithmetic model \(\mathcal X\) of a smooth projective curve \(X/K\), with dualizing sheaf \(\omega\) carrying the Arakelov metric, the source asks whether\[\begin{aligned}\widehat{\deg}(\omega^2)\stackrel{?}{=}&\ a_1(2g-2)\log|\Delta_{K/\Q}|+a_2\sum_{x\in\mathcal X_{\mathrm{sing}}}\log\#k(x)\\&+\sum_{\sigma:K\hookrightarrow\C}\delta_{\mathrm{Fal}}(X_\sigma(\C))+a_3\end{aligned}\] for constants \(a_1,a_2,a_3\in\R\). Here \(K\) is the intended number field, \(\Delta_{K/\Q}\) its absolute discriminant, \(g\) the genus of \(X\), \(k(x)\) the residue field at a closed arithmetic point, and \(\delta_{\mathrm{Fal}}\) the Faltings delta-invariant of the indicated complex fiber, with a convention still to be fixed. The archimedean sum is over the embeddings displayed in the source; no division by \([K:\Q]\) or weighting of conjugate pairs is inserted.\par A usual arithmetic model is a proper flat regular arithmetic surface over \(\operatorname{Spec}\mathcal O_K\) with generic fiber \(X\). If \(\bar\omega\) denotes the metrized dualizing line bundle, the left side is its arithmetic self-intersection \(\widehat{\deg}(\widehat c_1(\bar\omega)^2)\), not merely the degree of its tensor square on the generic curve. The source does not distinguish relative dualizing over \(\mathcal O_K\) from another dualizing convention, specify minimality or semistability of the model, or give the genus restrictions and metric normalization. These affect the proposed identity.\par The index set \(\mathcal X_{\mathrm{sing}}\) is essential but undefined. A regular scheme has no intrinsic nonregular points, so the intended set may instead be singular points of fibers or the nonsmooth locus over the arithmetic base. These are different notions. Without suitable fiber hypotheses, the latter need not be a finite set of points to which the displayed finite logarithmic sum applies. Multiplicities, local discriminant contributions, and any semistable-node convention are not supplied and cannot be reconstructed from the symbol alone.\par The dependence allowed for \(a_1,a_2,a_3\) is also unspecified. The nontrivial question would require constants uniform over a stated class, perhaps with stated dependence on genus or field. If \(a_3\) may be selected separately for each curve and model, the equality imposes no restriction. The source provides neither fixed values nor those uniformity quantifiers. The equality is preserved exactly rather than changed to an upper or lower bound, although the triage checked 2026-08-21 describes \emph{height and arithmetic-intersection bounds} as conjectural expectations and says no general proof was located. That mismatch and its dated open label require review; no established universal formula or current literature assessment is asserted here.
\subsection{Short English statement}
The source proposes an exact arithmetic self-intersection identity involving the field discriminant, bad-fiber contributions, and archimedean delta-invariants. Its constants' uniformity, dualizing and metric conventions, and singular-point sum are missing. The equality must not be replaced silently by the bounds mentioned in the attached assessment.
\subsection{Sources}
\begin{itemize}\item Local UnsolvedMath registry, permanent ID 30000326, code catalogue entry; \emph{Arakelov Self-Intersection Formula for Arithmetic Curves}, Oberwolfach Reports Open Problems.\item \emph{Arakelov Geometry (2006)}, Oberwolfach Reports, DOI 10.4171/owr/2005/43, pp. 2447--2492. \url{https://doi.org/10.4171/owr/2005/43}. Proposed year and year tag: 2005; no proposer named.\item Supplied literature label \emph{original Oberwolfach report}: \url{https://doi.org/10.4171/owr/2005/43}. No separate authoritative formula citation or normalization is supplied.\item Literature triage checked 2026-08-21, status \texttt{open}; record updated 2026-09-16. The \texttt{corrected\_verified} field claims only editorial and LaTeX cleanup, with no source-dependent content added. The dated assessment was not independently checked here.\end{itemize}
\end{document}
